
seulement si "@@M0@@ est sans torsion", i.e. @@M1@@ est isomorphe à un produit de groupes @@M2@@ et @@M3@@. Donc dans ce cas, et dans ce cas seulement, l'homomorphisme
@@M4@@
est un isomorphisme. Donc si @@M5@@ a de la torsion, le schéma sous-jacent fournit un exemple d'un @@M6@@ non bijectif.

\textbf{4.8.} On trouve de même un exemple où l'homomorphisme @@M7@@ de (4.1)
@@M8@@
ne se factorise pas par @@M9@@, ou ce qui revient au même ici, un exemple où cet homomorphisme n'est pas nul en degrés @@M10@@. Se limitant (comme il est loisible grâce à des théorèmes de spécialisation convenables) au cas où @@M11@@, on est ramené à étudier l'homomorphisme
@@M12@@
déduit d'ailleurs par passage au quotient de l'homomorphisme canonique
@@M13@@
Il est tentant de conjecturer que l'homomorphisme (4.10) est en fait injectif, puisqu'on sait par ailleurs que les nombres premiers de torsion pour @@M14@@ sont exactement ceux qui divisent l'ordre du groupe fini @@M15@@. S'il en était ainsi, dès que @@M16@@ aurait de la torsion, l'homomorphisme (4.10) ne se factoriserait pas par @@M17@@, donc pour @@M18@@ convenable (de façon précise, pour @@M19@@ divisant l'ordre de @@M20@@, i.e. @@M21@@ coefficient de torsion de @@M22@@) il en serait de même pour (4.9). On peut du moins affirmer que dans le cas du groupe @@M23@@ sur le corps @@M24@@, l'homomorphisme (4.10) n'est pas nul en degré @@M25@@, de façon précise il existe un élément de @@M26@@ dont l'image dans @@M27@@ est d'ordre @@M28@@; il en résulte donc que l'homomorphisme (4.9) ne se factorise pas par @@M29@@ lorsque @@M30@@. Le fait utilisé concernant @@M31@@, qui revient à dire qu'il existe dans @@M32@@ un élément dont l'image dans @@M33@@ est non nulle, se trouve dans [3], modulo le dictionnaire connu entre le langage groupes algébriques semi-simples complexes, et groupes de Lie compacts; il m'a été communiqué obligeamment par J.P. SERRE.

